\documentclass[11pt]{article}
\usepackage[margin=1in]{geometry}
\usepackage[T1]{fontenc}
\usepackage{lmodern,microtype}
\usepackage{amsmath,amssymb,amsthm,mathtools}
\usepackage{xcolor}
\usepackage[colorlinks=true,linkcolor=blue!40!black,citecolor=blue!40!black,urlcolor=blue!40!black]{hyperref}
\setlength{\emergencystretch}{2em}
\newcommand{\C}{\mathbb C}
\newcommand{\R}{\mathbb R}
\newcommand{\dd}{\,d}
\newcommand{\nr}{\mathrm{nr}}
\newcommand{\all}{\mathrm{all}}
\newcommand{\D}{\mathsf D}
\newcommand{\HH}{\mathcal H}
\newcommand{\FF}{\mathcal F}
\newcommand{\BB}{\mathcal B}
\newcommand{\NN}{\mathcal N}
\newcommand{\LL}{\mathcal L}
\newcommand{\WW}{\mathcal W}
\DeclareMathOperator{\supp}{supp}
\newtheorem{theorem}{Theorem}[section]
\newtheorem{proposition}[theorem]{Proposition}
\newtheorem{lemma}[theorem]{Lemma}
\newtheorem{corollary}[theorem]{Corollary}
\theoremstyle{remark}
\newtheorem{remark}[theorem]{Remark}
\hypersetup{pdftitle={Heat collisions and relative zero flux},pdfauthor={Drafted for Edward Baker}}
\title{Heat collisions and relative zero flux}
\author{Drafted for Edward Baker}
\date{Analytic manuscript\\23 October 2026}
\begin{document}
\maketitle

\begin{abstract}
A real collision in the backward heat equation consumes nonreal zeros.
We express this loss as an exact balance between an earlier population
and the flux through a moving window. If the window is defined by a
reference heat field $G$, its relative flux is governed by the single
coupled expression $H_{zz}G_z-H_zG_{zz}$: the apparent derivative
denominator in the boundary velocity cancels. We recover the flux from
smooth field integrals, prove sharp multiplicity bounds, and estimate the
reference mismatch without separating colliding roots. For the theta heat
flow, the mismatch has an exact antisymmetric integral involving the
unchanged arithmetic source and its complete tail. The ordered square
coefficients make the antisymmetrized source difference positive, giving
a sharper absolute estimate with quadratic diagonal cancellation. Local analytic
estimates make these bounds finite; obtaining useful uniform size at
growing height remains an arithmetic problem.
\end{abstract}

\section{The balance law}

Let $H$ be a nonzero, jointly holomorphic solution of
\begin{equation}
 H_t=-H_{zz},\qquad H(t,\bar z)=\overline{H(t,z)}\quad(t\in\R),
 \label{eq:heat}
\end{equation}
on a connected complex neighborhood of a compact time--space cylinder. We study real times
$I=[a,b]$. A simple zero moves with velocity
\begin{equation}
 v(t,z)=\frac{H_{zz}(t,z)}{H_z(t,z)}.
 \label{eq:velocity}
\end{equation}
Individual velocities become singular at a collision. The weighted
population nevertheless has finite variation.

A \emph{moving window} is a function $\psi(t,z)$, $C^1$ in time and $C^2$
in space, with $0\le\psi\le1$, $\psi(t,\bar z)=\psi(t,z)$, and support
contained in a fixed compact spatial set. Its derivatives are continuous
on the cylinder. Every zero in a sum below is counted with its full
spatial multiplicity. Put
\begin{align}
 N_{\psi,\nr}(t)&=\sum_{\substack{H(t,z)=0\\\Im z\ne0}}\psi(t,z),
 \label{eq:nonreal-population}\\
 \FF_{\psi,\nr}^{\pm}
 &=\int_a^b\sum_{\substack{H(t,z)=0\\\Im z\ne0}}
 \bigl[\pm\{\psi_t+\nabla\psi\cdot v\}\bigr]_+\dd t.
 \label{eq:relative-flux}
\end{align}
The integrands are defined at regular times; the exceptional times have
measure zero. Here $[s]_+=\max(s,0)$ and complex velocities are viewed as
planar vectors. These are one-sided fluxes: repeated inward and outward
motion is counted each time.

At a real multiple zero $(t_e,x_e)$ of order $m_e$, let
\begin{equation}
 \BB_\psi=\sum_{a<t_e\le b}
 2\left\lfloor\frac{m_e}{2}\right\rfloor\psi(t_e,x_e).
 \label{eq:birth-weight}
\end{equation}
The endpoint convention is deliberate: an event at $b$ is included; an
event at $a$ has already occurred.

\begin{theorem}[Nonreal mass and relative flux]
\label{thm:balance}
Under the preceding hypotheses, the fluxes are finite and
\begin{equation}
 \boxed{\quad
 N_{\psi,\nr}(b)+\BB_\psi
 =N_{\psi,\nr}(a)+\FF_{\psi,\nr}^{+}-\FF_{\psi,\nr}^{-}.
 \quad}
 \label{eq:balance}
\end{equation}
If $\psi_t+u\cdot\nabla\psi=0$ for a smooth boundary velocity $u$, then
the integrand in \eqref{eq:relative-flux} is
$[\pm\nabla\psi\cdot(v-u)]_+$.
\end{theorem}

Thus the relevant motion is relative to the boundary. Formula
\eqref{eq:balance} applies equally when the window changes shape without
being carried by a globally regular velocity field.

Suppose $\psi(b,x)=1$ at each real nonsimple endpoint site in a prescribed
core $K$. Let
\[
 M_K=\sum_{\substack{x\in K\cap\R\\m_x\ge2}}m_x,
 \qquad D_K=\sum_{\substack{x\in K\cap\R\\m_x\ge2}}(m_x-1),
 \qquad \mathcal R_\psi=N_{\psi,\nr}(a)+\FF_{\psi,\nr}^{+}.
\]

\begin{corollary}[Sharp endpoint bounds]
\label{cor:multiplicity}
\begin{equation}
 \boxed{D_K\le\mathcal R_\psi,\qquad
 M_K\le\frac32\mathcal R_\psi.}
 \label{eq:sharp-bounds}
\end{equation}
Both constants are optimal. If $\mathcal R_\psi<2$, there is no nonsimple
real endpoint site in $K$.
\end{corollary}

These statements concern real endpoint zeros and need no Riemann
hypothesis. Endpoint reality is an additional hypothesis when all zeros
of a canonical target are to be covered.

\subsection{Why a collision costs an even number of nonreal labels}

\begin{lemma}[Local heat birth]
\label{lem:birth}
At a spatial zero of order $m$ at $(t_e,z_e)$, the leading parabolic
polynomial, with $s=t-t_e$ and $w=z-z_e$, is a nonzero multiple of
\begin{equation}
 p_m(s,w)=\sum_{k=0}^{\lfloor m/2\rfloor}
 \frac{m!(-s)^k w^{m-2k}}{k!(m-2k)!}.
 \label{eq:hermite}
\end{equation}
The local roots are analytic in a one-sided square-root time parameter
and have integrable speeds $O(|s|^{-1/2})$. At a real event exactly
$2\lfloor m/2\rfloor$ nonreal labels disappear as real time increases.
Multiple events are isolated, and hence finite on the compact sets used
here.
\end{lemma}
\begin{proof}
Successive time derivatives are $(-\partial_z^2)^kH$, so the terms of
parabolic degree $m$ are precisely \eqref{eq:hermite}. For $s>0$,
$p_m(s,w)=(2s)^{m/2}\mathrm{He}_m(w/\sqrt{2s})$, where the monic
probabilists' Hermite polynomial has distinct real roots. For $s<0$ the
rescaled roots are purely imaginary, with one zero root if $m$ is odd.
The implicit function theorem after square-root rescaling continues each
simple rescaled root analytically. At a real event, real type and
uniqueness keep the forward roots real and the odd central backward
branch real. The other backward branches are nonreal.

The same rescaling is valid for complex $s$. It makes the local prepared
discriminant a nonzero multiple of
$s^{m(m-1)/2}$ to leading order. There is consequently no other multiple
event sufficiently near the given one. Compactness proves finiteness.
A spatial slice cannot vanish identically: the heat equation would make
all its time derivatives vanish, and analyticity would force $H$ to be
identically zero on the connected neighborhood.
\end{proof}

\begin{proof}[Proof of Theorem~\ref{thm:balance} and
Corollary~\ref{cor:multiplicity}]
Between events, differentiate the window along each root. Lemma
\ref{lem:birth} makes the resulting variation integrable. Roots entering
or leaving the compact support cause no jump because the window vanishes
there. An off-real collision preserves the weighted nonreal mass; a real
collision removes exactly the weight in \eqref{eq:birth-weight}.
Adding these changes proves \eqref{eq:balance}.

For $m\ge2$,
\[
 m-1\le2\lfloor m/2\rfloor,
 \qquad m\le\tfrac32\,2\lfloor m/2\rfloor.
\]
Each target site is counted with window value one. Discarding the
nonnegative terms $N_{\psi,\nr}(b)$ and $\FF^-_{\psi,\nr}$ proves
\eqref{eq:sharp-bounds}. A target multiple site costs at least two,
giving the final assertion. The field $H=z^3-6tz$, with a window equal
to one on its local roots, has earlier nonreal mass two and endpoint
values $D_K=2$, $M_K=3$. It attains both constants.
\end{proof}

\subsection{A boundary on the heat scale}

The backward double $z^2-2t$ suggests the contracting tube
\[
 h(t)^2=h_0^2-2t,\qquad
 \psi(t,x+iy)=\eta(x)\zeta(y/h(t)),\qquad t\le0,
\]
with $h_0>0$, smooth compact cutoffs, and $\zeta$ even. Its velocity is
$u=-iy/h(t)^2$. The relative flux pays horizontal motion as well as
vertical motion relative to this tube.

For the hard strip $(-A,A)+i(-h(t),h(t))$, the quadratic pair remains
inside. The cubic pair of $z^3-6tz$ instead enters at
$t=-h_0^2/4$ if $a<-h_0^2/4$. At the upper crossing, the root velocity
is $-3/h$ and the wall velocity is $-1/h$, so the relative inward speed
is $-2/h$. The two entering labels pay the endpoint triple. A boundary
adapted to the double's scale therefore still needs the flux term.

More generally, for smooth moving domains with simple transverse wall
crossings, the hard-wall analogue of \eqref{eq:balance} replaces
$\FF^\pm$ by entry and exit counts according to the sign of
$n\cdot(v-u)$. Endpoint wall zeros and multiple wall events need separate
conventions. Smooth windows avoid those boundary-event conventions.

\section{A natural moving boundary and its cancellation}
\label{sec:reference}

For a real window write $\D\psi=\psi_x-i\psi_y$ and
\[
 A_\psi=\psi_t|H_z|^2+\Re\{\D\psi\,H_{zz}\overline{H_z}\}.
\]
This is the smooth numerator of its zero flux;
Section~\ref{sec:field} gives its exact field-integral meaning.

Choose a real-type analytic reference $G$. Let $\chi:\C\to[0,1]$ be a
conjugation-symmetric $C^2$ cutoff of compact support and put
\begin{equation}
 \psi(t,z)=\chi(G(t,z)).
 \label{eq:reference-window}
\end{equation}
We use this window on a bounded spatial domain on whose collar
$G(t,z)$ avoids $\supp\chi$ for every $t\in I$, and extend it by zero.
This is the compact reference-packet condition. It can include several
components at once; no component is discarded when level sets merge.

For a radial cutoff, the transition boundaries are lemniscates
$|G(t,z)|=r$. Wherever $G_z\ne0$, their normal motion is that of
$-G_t/G_z$. If $G$ solves the heat equation, they follow its heat
motion. At a critical level there need not be a regular material
velocity; the smooth window \eqref{eq:reference-window} and the formula
below remain well defined.

\begin{theorem}[Reference-level flux]
\label{thm:reference}
For any analytic reference $G$, the moving-window numerator is
\begin{equation}
 A_\psi=\Re\{\D\chi(G)
              (G_tH_z-G_zH_t)\overline{H_z}\}.
 \label{eq:determinant-general}
\end{equation}
If both fields solve the heat equation, it is
\begin{equation}
 \boxed{\quad A_\psi=\Re\{\D\chi(G)\,
       \WW(H,G)\overline{H_z}\},\qquad
 \WW(H,G)=H_{zz}G_z-H_zG_{zz}.\quad}
 \label{eq:matched-determinant}
\end{equation}
Neither expression divides by $H_z$ or $G_z$.
\end{theorem}
\begin{proof}
The chain rule gives $\D_z\psi=\D\chi(G)G_z$ and
$\psi_t=\Re\{\D\chi(G)G_t\}$. Substitution into
\eqref{eq:kernel} gives \eqref{eq:determinant-general}; the heat equations
give \eqref{eq:matched-determinant}.
\end{proof}

The apparent singularity of the reference boundary velocity has canceled
against the gradient of its level window. If $G=cH$ with constant $c$,
the relative flux vanishes exactly. This is the material case, but it
does not estimate the earlier nonreal population.

\paragraph{Localizing without a reference-packet assumption.}
If the reference-packet condition is unavailable, take a fixed compact
spatial cutoff $\eta$ and use $\psi=\eta\chi(G)$. Then the exact
numerator is
\begin{equation}
 A_\psi=\eta\Re\{\D\chi(G)(G_tH_z-G_zH_t)\overline{H_z}\}
 +\chi(G)\Re\{\D\eta\,H_{zz}\overline{H_z}\}.
 \label{eq:outer-cutoff}
\end{equation}
The second term is the outer spatial flux. The positive part is taken
after adding the two terms; bounding them separately is a permissible
overestimate. Thus a natural reference boundary does not conceal the
cost of restricting to a finite height band.

\subsection{A quantitative mismatch bound}

Fix a compact region containing the window and its derivatives, and put
\[
 \NN=\int_I\#\{z\text{ in this region}:H(t,z)=0\}\dd t,
 \qquad
 \LL=\int_I\sum_{H(t,z)=0\text{ in this region}}|v(t,z)|\dd t.
\]
Both are complete-population quantities; independent analytic upper
bounds are given in Section~\ref{sec:anchors}.

\begin{theorem}[Flux from reference error]
\label{thm:mismatch}
For the window \eqref{eq:reference-window}, let $c$ be any constant and
$e=G-cH$. If $|e_t|\le E_t$ and $|e_z|\le E_1$ on the region, then
\begin{equation}
 \FF_{\psi,\nr}^+\le\FF_{\psi,\all}^+
 \le\|\D\chi\|_\infty(E_t\NN+E_1\LL).
 \label{eq:mismatch-bound}
\end{equation}
For two heat fields one may take $E_t=E_2$, where $E_2\ge|e_{zz}|$.
For $\chi(w)=\chi_0(|w|^2/r^2)$ and $|e|\le E_0$, the factor
$\|\D\chi\|_\infty$ can be replaced by
\begin{equation}
 C_\chi(E_0)=\min\left\{\|\D\chi\|_\infty,
                   \frac{2\|\chi_0'\|_\infty E_0}{r^2}\right\}.
 \label{eq:radial-improvement}
\end{equation}
One may minimize the entire bound over $c$, using the same $c$ in all
error terms.
\end{theorem}
\begin{proof}
Along a true root, $G=e$ and
\[
 \frac{d}{dt}G(t,z(t))=e_t+e_zv.
\]
The chain rule, positive-part inequality, and integration prove
\eqref{eq:mismatch-bound}. The radial refinement follows from
$|\D\chi(G)|=2|\chi_0'|\,|e|/r^2$ at the root.
\end{proof}

The estimate preserves the common motion before taking absolute values.
It needs neither a true-field nor a reference derivative floor. If the
error lies strictly inside the plateau of $\chi$ for the whole time
interval, the matched flux is actually zero. With an outer cutoff,
its separate flux in \eqref{eq:outer-cutoff} still has to be paid.

Endpoint coverage also has a source test. If $\chi=1$ on $|w|\le r_*$,
then $|G(b,z)-cH(b,z)|<r_*$ on a prescribed core inside the reference-packet domain makes the window equal
to one at every true endpoint zero there. This does not require knowing
those zeros in advance. For an outer cutoff additionally require
$\eta=1$ on that core.

\subsection{Removing a known carrier}

A high-height normalization may be useful, but its transport must be
retained. Let $B$ be holomorphic, nonvanishing, and of real type, $f=H/B$, $g=G/B$, and
\[
 \mathcal T_B=\partial_z^2+2\lambda\partial_z+c_B,
 \qquad \lambda=B_z/B,\qquad c_B=B_t/B+B_{zz}/B.
\]
Then $-f_t=\mathcal T_Bf$, and the current formulas use $-f_t$ in place
of $H_{zz}$. For two physical heat fields and $\psi=\chi(g)$, the
reference determinant becomes
\begin{equation}
 \boxed{\WW_B(f,g)=g_z\mathcal T_Bf-f_z\mathcal T_Bg
 =f_{zz}g_z-f_zg_{zz}+c_B(fg_z-gf_z).}
 \label{eq:carrier}
\end{equation}
The common drift cancels exactly. The zeroth-order term remains,
including at a finite field regulator. With $f=g+e$,
\[
 \WW_B=e_{zz}g_z-e_zg_{zz}+c_B(eg_z-ge_z).
\]
If instead $g_t+\mathcal T_Bg=d_g$, add $d_gf_z$ to
\eqref{eq:carrier}. The mismatch estimate still holds with the actual
$E_t$; in normalized variables $E_t=E_2$ is generally false. For two normalized
heat fields, $e=g-cf$ instead satisfies
\[
 E_t\le E_2+2\|\lambda\|_\infty E_1+\|c_B\|_\infty E_0.
\]

\section{Arithmetic information from theta and xi}
\label{sec:arithmetic}

For the canonical field we use
\begin{align}
 H(t,z)&=\int_0^\infty e^{tu^2}\Phi(u)\cos(zu)\dd u,
 &8H(0,z)&=\xi(1/2+iz/2),\label{eq:theta}\\
 \Phi(u)&=\sum_{n\ge1}\Phi_n(u),
 &\Phi_n(u)&=(2\pi^2n^4e^{9u}-3\pi n^2e^{5u})e^{-\pi n^2e^{4u}}.
 \nonumber
\end{align}
These normalizations agree with \cite{polymath}. Every summand is
positive on $u\ge0$, and the complete source has super-exponential
decay. Positivity supplies coarse geometric bounds. The arithmetic
reference below uses the actual integer-square coefficients themselves.

\subsection{The complete source tail enters only through a determinant}

For a fixed integer $N\ge1$, put
\[
 A_N(u)=\sum_{n\le N}\Phi_n(u),\qquad T_N(u)=\sum_{n>N}\Phi_n(u),
 \qquad G_N(t,z)=\int_0^\infty e^{tu^2}A_N(u)\cos(zu)\dd u.
\]
The reference $G_N$ is an exact heat field. The target remains the
unchanged $H$, including the entire tail $E_N=H-G_N$.

\begin{proposition}[Arithmetic head--tail identity]
\label{prop:arithmetic}
The reference-level residual is
\begin{align}
 \WW(H,G_N)
 &=\iint_{(0,\infty)^2}e^{t(u^2+v^2)}
            A_N(u)T_N(v)K_z(u,v)\dd u\dd v,\label{eq:head-tail}\\
 K_z(u,v)&=uv\{v\sin(zu)\cos(zv)-u\cos(zu)\sin(zv)\}\nonumber\\
 &=\frac{uv}{2}\{(v-u)\sin(z(u+v))
                     +(u+v)\sin(z(u-v))\}.\label{eq:kernel-cancel}
\end{align}
In particular $K_z(v,u)=-K_z(u,v)$ and $K_z(u,u)=0$.
\end{proposition}
\begin{proof}
The head--head terms cancel in
$H_{zz}(G_N)_z-H_z(G_N)_{zz}$, leaving
$(E_N)_{zz}(G_N)_z-(E_N)_z(G_N)_{zz}$.
Differentiate the complete integrals and multiply; Fubini is justified
by super-exponential decay. This gives the first kernel. The addition
formulas for sine give the second.
\end{proof}

Thus the boundary estimate sees a coupled head--tail interaction rather
than two separately large derivative products. More generally, for two
source densities $\Phi,\Psi$, with $H=H_\Phi$ and $G=H_\Psi$,
the same kernel gives the exact antisymmetrization
\[
 \WW(H_\Phi,H_\Psi)=\frac12\iint e^{t(u^2+v^2)}K_z(u,v)
 [\Psi(u)\Phi(v)-\Psi(v)\Phi(u)]\dd u\dd v.
\] A constant rescaling of the source
therefore makes the determinant vanish, just as the geometry requires.

For theta, the ordering of the actual coefficients sharpens this
antisymmetrization further.
\begin{proposition}[Ordered arithmetic residual]
\label{prop:ordered-source}
The ratio $T_N(u)/A_N(u)$ is strictly decreasing for $u\ge0$.
Consequently, with
\[
 D_N(u,v)=A_N(u)T_N(v)-A_N(v)T_N(u),
\]
one has $D_N(u,v)>0$ for $u>v$, and
\begin{equation}
 \boxed{\WW(H,G_N)=\int_{u>v>0}e^{t(u^2+v^2)}
                  D_N(u,v)K_z(u,v)\dd u\dd v.}
 \label{eq:ordered-source}
\end{equation}
\end{proposition}
\begin{proof}
Put $X=e^{4u}$. For $n>k\ge1$, the exact coefficients give
\[
 \frac{\Phi_n(u)}{\Phi_k(u)}
 =\frac{n^2}{k^2}\frac{2\pi n^2X-3}{2\pi k^2X-3}
                e^{-\pi(n^2-k^2)X},
\]
and hence
\[
 \partial_X\log\frac{\Phi_n}{\Phi_k}
 =-\pi(n^2-k^2)
  -\frac{6\pi(n^2-k^2)}{(2\pi n^2X-3)(2\pi k^2X-3)}<0.
\]
Thus $\Phi_k(u)\Phi_n(v)-\Phi_k(v)\Phi_n(u)>0$ for $u>v$.
Sum these positive differences over $k\le N<n$; absolute convergence
gives $D_N(u,v)>0$, equivalently the asserted ratio ordering. Split
\eqref{eq:head-tail} into its two half-planes and use the antisymmetry
of $K_z$ to obtain \eqref{eq:ordered-source}.
\end{proof}

There is a concrete improvement even in an absolute estimate. For
$|\Im z|\le b_0$, \eqref{eq:kernel-cancel} gives
\begin{align}
 |K_z(u,v)|\le\min\biggl\{&uv(u+v)e^{b_0(u+v)},\notag\\
 &\frac{uv|u-v|}{2}
 \left[e^{b_0(u+v)}+|z|(u+v)e^{b_0|u-v|}\right]\biggr\}.
 \label{eq:kernel-majorant}
\end{align}
The first entry gives the separate absolute source-moment bound; the second
retains the diagonal cancellation. Taking their minimum is never worse
and is strictly better on a neighborhood of $u=v>0$. This improvement
is relative to that specified source-moment majorant, not to bounds
that retain actual oscillatory values of the reference derivatives. The signed integral
\eqref{eq:head-tail} may be considerably smaller; no optimal arithmetic
estimate is asserted by taking absolute values.

Denote the symmetric majorant in \eqref{eq:kernel-majorant} by
$M_z(u,v)$. The ordered source gives the further bound
\[
 |\WW(H,G_N)|\le\int_{u>v>0}e^{t(u^2+v^2)}D_N(u,v)M_z(u,v)\dd u\dd v.
\]
For every fixed real $t$ this is strictly smaller than
$\iint e^{t(u^2+v^2)}A_N(u)T_N(v)M_z(u,v)\dd u\dd v$:
their difference is the positive integral with density
$2e^{t(u^2+v^2)}A_N(v)T_N(u)M_z(u,v)$ on $u>v>0$. This uses the actual ordered
square coefficients, not positivity alone. The remaining kernel is
oscillatory, so the source sign does not supply a sign for the flux.
Both $D_N$ and $K_z$ vanish on the diagonal: the ordered integrand is
locally $O((u-v)^2)$. The ratio ordering also survives positive
reweightings of the same square terms; it does not by itself distinguish
theta from weighted-square collision controls. No uniform factor of
improvement is claimed.

\subsection{An explicit tail allowance}

Let $m=N+1$ and
\begin{equation}
 q_m(u)=(1+1/m)^4e^{-\pi(2m+1)e^{4u}},\qquad
 \mathcal T_N(u)=
 \frac{2\pi^2m^4e^{9u}e^{-\pi m^2e^{4u}}}{1-q_m(u)}.
 \label{eq:tail-envelope}
\end{equation}
Then $0<q_m(u)<1$ and $0<T_N(u)\le\mathcal T_N(u)$. Indeed, after
discarding the negative term of $\Phi_n$, consecutive terms
$n^4e^{-\pi n^2e^{4u}}$ have ratio at most $q_m(u)$ for $n\ge m$.
The geometric sum proves the claim. Hence, on $-\epsilon\le t\le0$
and $|\Im z|\le b_0$,
\begin{equation}
 |\partial_z^j E_N(t,z)|\le
 \mathcal E_{j,N}(b_0):=
       \int_0^\infty u^j e^{b_0u}\mathcal T_N(u)\dd u.
 \label{eq:tail-jets}
\end{equation}
All coefficients and both infinite tails are retained in this bound.
The exact signed source may be used instead of its envelope.

Combining Theorem~\ref{thm:mismatch} with independent upper bounds
$\NN\le\epsilon n_*$ and $\LL\le L_*$ gives
\begin{equation}
 \boxed{\FF_{\psi,\nr}^{+}
 \le C_\chi(\mathcal E_{0,N})
       \{\epsilon n_*\mathcal E_{2,N}+L_*\mathcal E_{1,N}\}.}
 \label{eq:arithmetic-flux-bound}
\end{equation}
For a nonradial cutoff use $\|\D\chi\|_\infty$ in place of
$C_\chi$. The compact reference-packet condition, endpoint coverage,
and any outer-cutoff flux remain part of this application.

For every fixed compact band the tail budgets tend to zero as $N$
increases. This does not imply a useful normalized estimate at growing
height. In particular, the reference's finite source prefix does not
inherit the full source reflection, and small absolute source errors
can dominate an exponentially small high-height field. The unchanged
arithmetic in \eqref{eq:head-tail} is available for a stronger estimate;
positivity and the envelope alone do not settle that estimate.

\subsection{A local reference written directly in xi}

An alternative uses the exact initial field $H_0(z)=\xi(1/2+iz/2)/8$:
\begin{equation}
 G_K(t,z)=\sum_{k=0}^{K}\frac{(-t)^k}{k!}H_0^{(2k)}(z),
 \qquad
 H_0^{(j)}(z)=\frac18(i/2)^j\xi^{(j)}(1/2+iz/2).
 \label{eq:xi-reference}
\end{equation}
It is an analytic reference with a known heat defect,
\begin{equation}
 d_K=(G_K)_t+(G_K)_{zz}
       =\frac{(-t)^K}{K!}H_0^{(2K+2)}.
 \label{eq:xi-defect}
\end{equation}
Its relative flux uses
$\WW(H,G_K)+d_KH_z$, as follows directly from
\eqref{eq:determinant-general}. Equivalently, with $e=G_K-H$,
\[
 (G_K)_zH_{zz}+(G_K)_tH_z=e_zH_{zz}+e_tH_z.
\]
The temporal-error form includes the defect exactly, so it is paid only
once in Theorem~\ref{thm:mismatch}.

For $A_j(b_0)=\int_0^\infty u^je^{b_0u}\Phi(u)\dd u$, the real
negative-time Taylor remainder gives, for $K\ge1$ and $e=G_K-H$,
\begin{equation}
 |\partial_z^j e|\le\frac{\epsilon^{K+1}}{(K+1)!}A_{j+2K+2}(b_0),
 \qquad
 |e_t|\le\frac{\epsilon^K}{K!}A_{2K+2}(b_0).
 \label{eq:xi-error}
\end{equation}
These inequalities follow under the complete integral from the remainder
of $e^{tu^2}$ for $t\le0$; they are not continuations of a
positive-time error theorem. The exact xi coefficients can be used in
the signed determinant before resorting to these absolute remainders.

\section{Recovering the flux from the complete field}
\label{sec:field}

Recall the moving-window numerator $A_\psi$ and define
\begin{equation}
 k_\rho(w)=\frac{\rho^2}{\pi(|w|^2+\rho^2)^2},\qquad
 A_\psi=\psi_t|H_z|^2+\Re\{\D\psi\,H_{zz}\overline{H_z}\}.
 \label{eq:kernel}
\end{equation}
The kernel has integral one on the target plane. Its pullback gives a
smooth density and current
\[
 \mu_\rho=k_\rho(H)|H_z|^2,\qquad
 j_\rho=k_\rho(H)H_{zz}\overline{H_z},\qquad
 \partial_t\mu_\rho+\operatorname{div}j_\rho=0.
\]
One can obtain the last identity by pulling back the closed target-plane
area form. The moving-window integrand is exactly
$k_\rho(H)A_\psi=\psi_t\mu_\rho+\nabla\psi\cdot j_\rho$. In a material window it is
$\nabla\psi\cdot(j_\rho-u\mu_\rho)$.

Let $0\le\vartheta_\delta(y)\le1$ be continuous, vanish for $|y|\le\delta$,
equal one for $|y|\ge2\delta$, and increase to $1_{y\ne0}$ as
$\delta\downarrow0$. Set
\begin{equation}
 Q_{\rho,\delta}^{\pm}
 =\int_I\int\vartheta_\delta(\Im z)k_\rho(H)
          [\pm A_\psi]_+\dd A_z\dd t.
 \label{eq:regularized-flux}
\end{equation}

\begin{theorem}[Exact field representation]
\label{thm:current-limit}
The complete nonreal flux is
\begin{equation}
 \FF_{\psi,\nr}^{\pm}
   =\lim_{\delta\downarrow0}\lim_{\rho\downarrow0}
                    Q_{\rho,\delta}^{\pm}.
 \label{eq:nonreal-limit}
\end{equation}
Without $\vartheta_\delta$, the limit is the corresponding variation of
\emph{all} roots. In particular it bounds the nonreal positive flux
from above. There is no excess current concentrated at the collision
times.
\end{theorem}

The measuring cutoff in \eqref{eq:regularized-flux} is not part of the
moving window. Replacing $\psi$ by $\psi\vartheta_\delta$ would introduce
additional derivative terms and would be a different formula.

\begin{lemma}[Uniform variation of nearby levels]
\label{lem:levels}
Let $f$ be jointly holomorphic near a compact cylinder, with no identically
zero spatial slice. On a compact spatial set, the roots of $f(t,z)=w$
have locally bounded total variation, uniformly for almost every
sufficiently small $w$. Their speed sums have a uniform $L^p(I)$ bound
for some $p>1$.
\end{lemma}
\begin{proof}
Cover the compact zero set by finitely many time--disk charts on whose
spatial circles $f$ stays away from zero, throughout a complex time
collar. Outside the chart interiors the remaining compact set has a
positive minimum of $|f|$. Choose the target radius below this minimum
and the circle bounds, so that every nearby level is covered. In each
chart $f-w$ has a fixed-degree monic root polynomial, with coefficients
holomorphic in time and $w$. After normalizing its disk and time interval,
root containment bounds its coefficients; Cauchy's estimates bound their
$C^{n-1,1}$ norms uniformly in $w$.

The polynomial-root theorem of Parusi\'nski--Rainer
\cite[Theorem 1]{pr} gives $L^p$ derivatives of every continuous root
for $1<p<n/(n-1)$. Choose one admissible $p$ for the finite cover and
sum the bounds. Degree one is immediate. Returning to physical time
introduces the factor $|J|^{1/p-1}$ on each chart; all these factors
and any overlaps are included. No root separation is used.
\end{proof}

\begin{proof}[Proof of Theorem~\ref{thm:current-limit}]
For almost every target $w$, the spatial area formula, whose Jacobian
is $|H_z|^2$, rewrites the integral without the measuring cutoff as
\begin{equation}
 Q_\rho^{\pm}=\int_\C k_\rho(w)V^{\pm}(w)\dd A_w,
 \quad
 V^{\pm}(w)=\int_I\sum_{H(t,z)=w}
 [\pm\{\psi_t+\nabla\psi\cdot v_w\}]_+\dd t,
 \label{eq:area}
\end{equation}
where $v_w=-H_t/H_z$. Quotients are used only at regular preimages;
critical points contribute zero to the original density. Exceptional
fibers are handled by the almost-everywhere area formula and Fubini.

At every regular time for the zero level, roots and velocities converge
as $w\to0$. The additional term $\psi_t$ is bounded by the finite chart
degrees. Lemma~\ref{lem:levels} makes the velocity terms uniformly
integrable across the finitely many zero-level events. Hence
$V^{\pm}(w)\to V^{\pm}(0)$ in the essential target-plane sense. The
portion $|w|\ge r$ in \eqref{eq:area} is at most
\[
 \frac{\rho^2}{\pi(r^2+\rho^2)^2}
 \int_I\int\bigl(|\psi_t||H_z|^2+
                   |\nabla\psi||H_tH_z|\bigr)\dd A\dd t,
\]
which tends to zero. The approximate identity therefore converges to
the all-root flux. For fixed $\delta$, insert the continuous bounded
factor $\vartheta_\delta$ in the same proof; then apply monotone
convergence on the root trajectories as $\delta\downarrow0$.
\end{proof}

The earlier population has a parallel field representation:
\begin{equation}
 N_{\psi,\nr}(a)=
 \lim_{\delta\downarrow0}\lim_{\rho\downarrow0}
 \int\vartheta_\delta(\Im z)\psi(a,z)
                   k_\rho(H(a,z))|H_z(a,z)|^2\dd A_z.
 \label{eq:initial-field}
\end{equation}
At a multiple zero the pullback area counts its full multiplicity.
Thus \eqref{eq:balance}, \eqref{eq:nonreal-limit}, and
\eqref{eq:initial-field} describe the whole nonreal resource through
the complete field. They do not yet make that resource small.

\begin{remark}[The near-real strip]
For a conjugation-symmetric moving window the omitted flux is locally
$O(\delta^2)$. Indeed, near a real event put $q=\sqrt{t_e-t}$.
Uniqueness of the rescaled branches and real type give
$z_j(-q)=\overline{z_j(q)}$. Therefore $x_j$ is even and $y_j$ odd
in $q$. In time derivatives, $x_j'=O(1)$,
$y_j'=O(q^{-1})$, while $\psi_y=O(y_j)=O(q)$.
The total window derivative is bounded. If
$|y_j|\ge c_j\sqrt{t_e-t}$ and $|d\psi/dt|\le A_j$, the omitted
portion costs at most $4A_j\delta^2/c_j^2$. Sum over the finite incoming
branches after taking $\delta$ below the other tracks' clearance.
These constants are local field data; no uniform height estimate is
implicit in this observation.
\end{remark}

\section{Independent bounds and the remaining arithmetic question}
\label{sec:anchors}

The geometric quantities in \eqref{eq:arithmetic-flux-bound} need not
be supplied by individual zero tracking.

\begin{theorem}[An anchored analytic bound]
\label{thm:anchor}
Let $f$ be holomorphic on a neighborhood of
\[
 \{\tau:\operatorname{dist}(\tau,I)\le d\}\times\{|z|\le4R\},
\]
with $|f|\le M$ and $|f(t,0)|\ge m>0$ for real $t\in I$.
Write $\alpha=m/M$. There are constants $n_*(\alpha)$ and
$C(\alpha)$ such that the roots in $|z|\le R$ obey
\begin{equation}
 \#\{f(t,z)=0:|z|\le R\}\le n_*(\alpha),\qquad
 \int_I\sum_{\substack{f(t,z)=0\\|z|\le R}}|z'(t)|\dd t
 \le C(\alpha)R(1+|I|/d).
 \label{eq:anchor-bound}
\end{equation}
No lower bound on $f_z$ is required. The constants depend only on
$\alpha$ and the fixed radius ratios; they are not numerically
evaluated here.
\end{theorem}
\begin{proof}
Jensen bounds the zero count in $|z|<3R$ by
\[
 n=\max\left(1,\left\lceil\frac{\log(M/m)}{\log(4/3)}\right\rceil\right).
\]
At a fixed real time remove these zeros with disk Blaschke factors for
the $3R$ disk. The remaining zero-free factor has modulus at most $M$
and modulus at least $m$ at zero. Writing this factor as $h$, Harnack applied to the nonnegative harmonic
function $\log(M/|h|)$ bounds its modulus below by $M\alpha^5$ on $|z|\le2R$.
Choose $s\in[R,2R]$ at distance at least $R/(4n)$ from all the removed
zero moduli. At most half the radius interval is excluded, so such a
circle exists, and on it
\[
 |f(t,z)|\ge M\alpha^5(20n)^{-n}=:\eta.
\]
Zeros on the factorization circle are only boundary zeros and do not
affect Harnack in the interior.

Cauchy's time estimate bounds $|f_t|$ by $2M/d$ on the half-width
tube. Put $\beta=\eta/M$ and $\ell=d\beta/16$, and partition $I$
into at most $1+|I|/\ell$ intervals of length at most $\ell$.
Each midpoint's circle stays clear on its interval and a complex-time
collar of width $\ell/2$: the field changes by at most
$2M\ell/d=\eta/8$. Each chart encloses at most $n$ roots.
Normalize its time interval and disk; the normalized time collar has
width at least $1/2$. The binomial coefficient bounds, Cauchy's estimates,
and the polynomial-root theorem used in Lemma~\ref{lem:levels} give
\[
 \int_J\sum|z'(t)|\dd t
 =s\int_0^1\sum|\zeta_j'(u)|\dd u\le C_2(n)R.
\]
The physical time scale cancels in this path-length integral. Summing proves \eqref{eq:anchor-bound}; Jensen
also gives the stated pointwise population bound.
\end{proof}

For $f=H$ on $I=[-\epsilon,0]$, actual complete-source data are
\begin{equation}
 m=\int_0^\infty e^{-\epsilon u^2}\Phi(u)\dd u,
 \qquad
 M=\int_0^\infty e^{du^2+4Ru}\Phi(u)\dd u.
 \label{eq:theta-anchor}
\end{equation}
They are finite and positive, and they meet the theorem's hypotheses.
They give a coarse earlier nonreal bound as well as the root-path bound
needed in \eqref{eq:arithmetic-flux-bound}. A distant band can be
enclosed in such a disk, but the resulting constants can be very large.

Arithmetic also supplies anchors away from the critical line. For
$z_*=h+ib_0$, $b_0>1$, put $s_*=(1+b_0)/2-ih/2$ and
$\sigma=\Re s_*>1$. The functional equation \cite[\S25.4]{dlmf} and Euler product give
\begin{equation}
 |H(0,z_*)|\ge
 \frac{|s_*(s_*-1)\Gamma(s_*/2)|}{16\pi^{\sigma/2}}
 \frac{\zeta(2\sigma)}{\zeta(\sigma)}=:m_0.
 \label{eq:euler-anchor}
\end{equation}
Here $|\zeta(\sigma+it)|\ge\prod_p(1+p^{-\sigma})^{-1}$.
For negative real time,
$|H(t,z_*)-H(0,z_*)|\le\epsilon A_2(b_0)$.
Thus $m_0-\epsilon A_2(b_0)>0$ is a genuine local anchor when that
inequality holds. It can be used on local disks, and reflected disks
when symmetry is needed. This elementary continuation may force an
extremely small width at large height; no favorable microscopic width
is inferred from it.

\paragraph{What remains to estimate.}
The central formula is exact, and its multiplicity constants are sharp.
The reference determinant removes a derivative denominator and retains
the arithmetic cancellation before estimation. The bounds above give
finite, independently sourced allowances. What is still needed for a
useful global conclusion is a sufficiently small earlier \emph{nonreal}
population together with sufficiently small complete relative flux on
prescribed height bands, including outer localization and coverage.
The generic anchor bound does not discriminate real from nonreal zeros.
The arithmetic tail bounds and signed determinant have no proved useful
height-uniform size here. Neither RH nor a new global simplicity
proportion follows from the present results.

For a fixed moving-window family, extending its initial time backward
cannot decrease its exact conservative resource: rearranging
\eqref{eq:balance} gives
$\mathcal R_\psi=N_{\psi,\nr}(b)+\BB_\psi+\FF^-_{\psi,\nr}$.
A better completed estimate must come from better arithmetic information
or a better justified window, with the same endpoint coverage.

\appendix
\section{Small values and finite-regulator stability}

The proof of Lemma~\ref{lem:levels} gives more than finiteness. For a
fixed analytic $f$ and moving window put
\[
 W_\psi(r)=\int_I\int_{|f|\le r}
       \bigl(|\psi_t||f_z|^2+|\nabla\psi||f_tf_z|\bigr)\dd A\dd t.
\]
The area formula and the chart bounds give $W_\psi(r)\le\pi Lr^2$
for small $r$, with $L$ bounded by the corresponding population and
path-length data. The target-plane coefficient of $r^2$ has a finite
limit when the zero-level multiple events are finite. Integrating
$k_\rho$ against this distribution uniformly bounds the regularized
currents. Its limit is the corresponding absolute-variation bound; the
exact directional limit retains the coupled numerator as in
Theorem~\ref{thm:current-limit}.

That limit has no uniform rate under degeneration of the field or its
charts. At a fixed regulator there is, however, a direct error estimate.
For any analytic $f$ use
\[
 A_f=\psi_t|f_z|^2+\Re\{\D\psi(-f_t)\overline{f_z}\}.
\]
If $g$ approximates $f$, set $e_0\ge|f-g|$ and
$e_A\ge|A_f-A_g|$. The elementary kernel bounds
\[
 \|k_\rho\|_\infty\le\frac1{\pi\rho^2},\qquad
 \|\nabla k_\rho\|_\infty\le\frac2{\pi\rho^3}
\]
and the Lipschitz property of positive part give
\[
 |Q_\rho(f)-Q_\rho(g)|
 \le\int_I\int\left\{\frac{e_A}{\pi\rho^2}
               +\frac{2e_0|A_g|}{\pi\rho^3}\right\}\dd A\dd t.
\]
Subtract the coupled numerator before bounding $e_A$. This statement
holds for a fixed window; an approximation-dependent change of window
needs its own error allowance. A fixed source error cannot be passed
unchanged through $\rho\downarrow0$.

\section*{Acknowledgements and status}
This manuscript was drafted for Edward Baker with substantial assistance
from OpenAI Codex, model \texttt{gpt-6-astra}, reasoning effort
\emph{ultra}. Contributing and reviewing agents used the same model.
The finite analytic statements have undergone internal derivation checks;
these are not independent specialist refereeing or formal verification.
The earlier research draft and accompanying notes retain the broader
coefficient, spectral, contour and numerical investigations. This version
selects the moving-boundary conservation law and its source estimates as
one continuous argument. No priority claim is made for the classical
ingredients or their combination.

\begin{thebibliography}{9}
\bibitem{pr}
A.~Parusi\'nski and A.~Rainer,
\emph{Optimal Sobolev regularity of roots of polynomials},
Annales scientifiques de l'\'Ecole normale sup\'erieure (4) 51 (2018),
1343--1387. \href{https://doi.org/10.24033/asens.2376}{doi:10.24033/asens.2376}.

\bibitem{polymath}
D.~H.~J. Polymath,
\emph{Effective approximation of heat flow evolution of the Riemann
$\xi$ function, and a new upper bound for the de Bruijn--Newman constant},
arXiv:1904.12438v2 (2019).
\url{https://arxiv.org/abs/1904.12438}.
\bibitem{dlmf}
NIST Digital Library of Mathematical Functions, \S25.4,
\emph{Reflection formulas}. \url{https://dlmf.nist.gov/25.4}.
\end{thebibliography}
\end{document}
